\verseline{}{众人静默下来，凝神望着他的面庞。}
\verseline{}{父亲埃涅阿斯便从高高的长榻上开口：}
\verseline{}{“王后，你命我重新唤起难以言说的悲痛：}
\verseline{}{讲达那俄斯人怎样摧毁特洛伊的强盛和那可悲的王国，}
\verseline{5}{讲我亲眼见过的至惨景象，}
\verseline{}{以及我深陷其中的一切。说起这些，}
\verseline{}{即使密耳弥冬人、多洛佩斯人或冷酷的乌利塞斯麾下士兵，\footnote{密耳弥冬人是阿喀琉斯的部众；多洛佩斯人也属其阵营。乌利塞斯即希腊史诗中的奥德修斯，拉丁文名为 Ulixes。埃涅阿斯说连这些冷酷敌人听到特洛伊覆亡也会流泪，是以敌方为尺度强调故事之惨。}}
\verseline{}{又有谁忍得住眼泪？湿润的黑夜正从天空}
\verseline{}{沉落，西坠的群星催人入眠。}
\verseline{10}{但既然你如此渴望知道我们的遭际}
\verseline{}{想简短听完特洛伊最后的苦难，}
\verseline{}{尽管我的心一想起便战栗，因哀痛而退缩，}
\verseline{}{我仍要开口。达那俄斯的诸将被战争击溃，}
\verseline{}{又为命运所拒。多年流逝以后，}
\verseline{15}{他们凭帕拉斯神圣的技艺造出山岳般的木马，}
\verseline{}{以刨削的冷杉木编结它的肋骨：\footnote{木马由弥涅耳瓦／帕拉斯的技艺启发而造。女神主管工艺与谋略，西农的骗局又把木马说成补偿帕拉狄乌姆被盗而献给她的祭物；这里的“神圣”首先同帕拉斯有关。}}
\verseline{}{假称这是为归航许下的还愿礼，传言便这样流布。}
\verseline{}{他们暗中抽签选出一批精锐勇士，}
\verseline{}{把他们关进木马幽暗的侧腹，让巨大的洞腔}
\verseline{20}{和腹部深处塞满全副武装的士兵。}
\verseline{}{视野之内有座声名远播的忒涅多斯岛\footnote{忒涅多斯是特洛伊近海岛屿。希腊舰队假作返乡，实藏在岛的荒岸边，稍后从这里驶回；西农举火示意的另一叙事传统见注73。}}
\verseline{}{普里阿摩斯王国尚存时，它财富丰饶，}
\verseline{}{如今只剩一个海湾和一处并不可靠的泊地：}
\verseline{}{他们驶到那里，藏在荒凉的岸上。}
\verseline{25}{我们以为他们已经离去，乘风驶回迈锡尼：}
\verseline{}{于是整个特洛伊终于从漫长的哀痛中解脱。}
\verseline{}{城门敞开，人们欣然走出观看多利安人的营地}
\verseline{}{观看荒废的阵地和被遗弃的海岸。}
\verseline{}{这里曾驻扎多洛佩斯人，这里凶猛的阿喀琉斯曾扎营；}
\verseline{30}{这里是舰队的泊处，也是他们列阵交锋的地方。\footnote{阿喀琉斯、密耳弥冬人和多洛佩斯人来自《伊利亚特》的希腊阵营。特洛伊人游览被弃营地，既像幸存者重访战场，也因把敌军撤离当成终战而处在全卷最危险的误读中。}}
\verseline{}{一些人惊望处女弥涅耳瓦那致命的礼物，}
\verseline{}{赞叹木马的庞大。提摩特斯第一个}
\verseline{}{劝人把它拖进城墙，安置在城堡上，}
\verseline{}{不知出于背叛，还是特洛伊的命数已这样催逼。}
\verseline{35}{但卡皮斯和那些见识更明的人，}
\verseline{}{主张把这藏有达那俄斯人诡计的可疑礼物}
\verseline{}{推入海中，或在下面放火焚烧，}
\verseline{}{或钻穿它空洞的腹部，探明藏匿之处。}
\verseline{}{犹疑的人群分裂成相反的两派。}
\verseline{40}{这时拉奥孔走在所有人前面，一大群人跟随，}
\verseline{}{他像一团怒火从城堡最高处奔下，}
\verseline{}{远远喊道：‘不幸的同胞们，何至于这样疯狂？}
\verseline{}{你们真信敌人已经走了？真以为达那俄斯人的礼物}
\verseline{}{会不带诡计？这就是你们所认识的乌利塞斯？}
\verseline{45}{要么阿凯亚人藏在这木头之中，}
\verseline{}{要么这器械是为攻打我们的城墙而造，}
\verseline{}{好窥探我们的家宅，再从高处压向全城，}
\verseline{}{或另有诡谋藏在其中。特乌克罗斯人，别信这匹马。\footnote{拉奥孔名言的拉丁原句为 timeo Danaos et dona ferentis，意为“我怕达那俄斯人，即使他们带着礼物”。重点不只是怀疑礼物，而是十年战争之后，即使达那俄斯人带着礼物，他们仍是敌人。}}
\verseline{}{无论它是什么，我都怕达那俄斯人，哪怕他们带着礼物来。’}
\verseline{50}{说罢他奋起全身力量，将一支巨矛}
\verseline{}{掷向木兽的侧腹，掷向那由木板弯合的腹部。}
\verseline{}{长矛插入兀自颤动，木马腹腔为之一震}
\verseline{}{深处的空洞轰然作响，发出呻吟。}
\verseline{}{倘若神意未曾与我们为敌，倘若我们的心不曾昏乱，}
\verseline{55}{他本会驱使我们以铁器剖开阿尔戈斯人的藏身处，}
\verseline{}{特洛伊如今本会依然矗立；普里阿摩斯巍峨的城堡啊，你也会依旧存在。}
\verseline{}{看，就在这时达耳达诺斯的牧人高声呼喊}
\verseline{}{把一个双手反缚在背后的青年拖到国王面前。}
\verseline{}{这陌生人主动迎向他们，把自己交出，}
\verseline{60}{正为完成这场诡计，替阿凯亚人打开特洛伊，}
\verseline{}{他心中无惧，早已准备好两种结局：}
\verseline{}{或施展骗局，或迎向注定的死亡。}
\verseline{}{特洛伊青年急于观看，从四面八方}
\verseline{}{涌聚过来争相戏弄俘虏。}
\verseline{65}{如今听听达那俄斯人的奸计吧，从这一桩罪恶}
\verseline{}{便可认清他们全部。}
\verseline{}{他惶恐地站在众目中央，手无寸铁}
\verseline{}{环视佛里吉亚人的队列，开口说道：}
\verseline{}{‘唉，如今还有哪片土地、哪片海洋}
\verseline{70}{肯收留我？我这可怜人最后还剩下什么？}
\verseline{}{达那俄斯人中已无我容身之处，连达耳达诺斯人}
\verseline{}{也满怀敌意，要以我的鲜血索取惩罚？’}
\verseline{}{这一声悲叹扭转了人心，众人的冲动全都平息。}
\verseline{}{我们催他开口，说出自己出自何种血脉，}
\verseline{75}{带来什么，也说明身为俘虏有何指望。\footnote{Conte 2019 正文由第75行直接接第77行。该版校勘记注明传统第76行与3.612相同，并在多个主要传本中缺失。}}
\verseline{}{‘国王，不论结果如何，我都会向你说出}
\verseline{}{一切真话，也绝不否认自己出身阿尔戈斯。}
\verseline{}{先承认这一点。残酷的命运纵使使西农不幸\footnote{西农是希腊人为使木马入城而留下的欺骗者。马克罗比乌斯（6.1.57）把他这里“命运虽能使我不幸，却不能使我成为骗子”的自辩与阿克齐乌斯《忒勒福斯》中“命运能夺走王国与财富，却不能夺走德性”的句式并列。维吉尔让真正的骗子借悲剧人物捍卫德性的语言为自己担保，因此这句“真话”从一开始就带着反讽；后文长篇供述同样是西农精心安排的骗词。}}
\verseline{80}{也不会把他变成虚妄的骗子。}
\verseline{}{倘若帕拉墨得斯——贝卢斯的后裔——这个名字}
\verseline{}{以及他显赫的声誉曾在传闻中抵达你的耳边：\footnote{帕拉墨得斯是希腊联军中的智者，后被乌利塞斯构陷处死。西农借这桩足以激起同情的希腊内部冤案，把自己也说成受乌利塞斯迫害的人，使谎话更容易取信。}}
\verseline{}{佩拉斯吉人以虚假的叛国罪名，}
\verseline{}{用可耻的指控把无辜的他送上死路，只因他反对战争，}
\verseline{85}{如今他失去生命，他们却又哀悼他。}
\verseline{}{他是我的亲族，我贫穷的父亲}
\verseline{}{从我年少时便派我随他参战。}
\verseline{}{当他仍安然在位，在诸王的议事中}
\verseline{}{声望正盛，我也曾有过一点名位与荣誉。}
\verseline{90}{后来，由于狡诈的乌利塞斯心怀嫉妒，}
\verseline{}{——我说的并非无人知晓——帕拉墨得斯离开了人世，}
\verseline{}{我受尽打击，在黑暗与哀痛中拖着生命，}
\verseline{}{为无辜朋友的命运愤愤不平。}
\verseline{}{我这疯子没有沉默，还曾发誓只要命运允许，}
\verseline{95}{只要有一天能以胜利者身份回到祖国阿尔戈斯，}
\verseline{}{我一定替他复仇。这些话激起了尖锐的仇恨。}
\verseline{}{祸根便从这里开始。从此乌利塞斯}
\verseline{}{不断用新的罪名恐吓我，在人群中散布}
\verseline{}{含混的流言，暗中寻找置我于死地的手段。}
\verseline{100}{他始终不肯罢休，直到借卡尔卡斯之手——}
\verseline{}{可我为什么还要徒然翻出这些令人厌恶的旧事？}
\verseline{}{倘若你们把所有阿凯亚人都看作一类，}
\verseline{}{听到这一点便够了，我又何必拖延？现在就惩罚我吧，}
\verseline{}{伊塔卡人正求之不得，阿特柔斯二子也愿重金买下。’}
\verseline{105}{于是我们反倒急切追问，想弄清其中缘由，}
\verseline{}{全然不识这样深重的罪恶和佩拉斯吉人的手段。}
\verseline{}{他装作惊惧，继续编造：}
\verseline{}{‘达那俄斯人屡次想抛下特洛伊}
\verseline{}{逃离这场漫长而令人疲惫的战争；}
\verseline{110}{但愿他们真做到了！海上的恶劣风暴一次次}
\verseline{}{封住归路，南风也惊退即将启航的人。}
\verseline{}{尤其当这匹枫木梁架成的木马}
\verseline{}{已经立在这里，整个天空都响起雷雨的轰鸣。}
\verseline{}{我们惶惶不安，派欧律皮罗斯去求问福玻斯的神谕，}
\verseline{115}{他从神殿深处带回这些悲苦的话：}
\verseline{}{“达那俄斯人，你们最初来到伊利乌姆海岸时，\footnote{西农所说的“被杀少女”通常指伊菲革涅亚：希腊舰队出征前在奥利斯受阻，阿伽门农之女被要求作为祭品。古代传统对她的结局有不同版本；西农把出征献祭套用到归航。}}
\verseline{}{曾以鲜血、以一位被杀的少女平息风暴；}
\verseline{}{如今归途也要用鲜血求得，要献上一条\footnote{卡尔卡斯是希腊先知。西农声称乌利塞斯强迫他指定人祭，故事把先知权威、集体恐惧与预先串通的阴谋拼成一套看似合乎祭仪的骗局。}}
\verseline{}{阿尔戈斯人的性命。”这声音一传入众人耳中，}
\verseline{120}{众人怔住，冰冷的战栗深入}
\verseline{}{骨髓：命运为谁准备？阿波罗索取的是谁？}
\verseline{}{这时伊塔卡人在巨大的喧嚷中把先知卡尔卡斯}
\verseline{}{拖到人群中央，逼问众神的旨意究竟是什么。}
\verseline{}{许多人早已向我预言那阴谋家的残酷罪行，}
\verseline{125}{又默默看出了将要发生的事。}
\verseline{}{卡尔卡斯整整十天沉默，藏起自己，不肯}
\verseline{}{亲口出卖任何人，把任何人推向死亡。}
\verseline{}{终于他被伊塔卡人巨大的叫嚷逼迫，}
\verseline{}{依照预先串通的安排开口，把我指定为祭坛上的牺牲。}
\verseline{130}{所有人都赞同；他们眼见各自所怕的灾祸，}
\verseline{}{转成我这可怜人一人的毁灭，便都坦然接受。}
\verseline{}{可怕的日子已经临近，祭仪为我备好，}
\verseline{}{咸麦粉和缠绕太阳穴的圣带也已备好。}
\verseline{}{我承认，我从死亡中挣脱，扯断绳索，}
\verseline{135}{趁夜藏进泥泞池沼的芦苇}
\verseline{}{等待他们扬帆——倘若他们果真扬帆。}
\verseline{}{如今我已不再指望看见古老的祖国，}
\verseline{}{也见不到亲爱的孩子与渴望重逢的父亲；}
\verseline{}{也许他们会因我的逃脱，向那些可怜人}
\verseline{140}{索取惩罚，以他们的死为我的过错赎罪。}
\verseline{}{因此我以天上诸神、以知悉真相的神明求你，}
\verseline{}{也以凡人世间若还留有任何}
\verseline{}{未受玷污的信义求你：怜悯这样深重的苦难，}
\verseline{}{怜悯这颗承受不应承受之事的心。’}
\verseline{145}{他的眼泪为他换得性命，我们甚至主动怜悯他。}
\verseline{}{普里阿摩斯首先命人解开他手上的镣铐}
\verseline{}{和紧绳，又用友善的话对他说：}
\verseline{}{‘无论你是谁，从此忘掉那些已经离去的希腊人；}
\verseline{}{你将成为我们的人。现在如实回答我的询问：}
\verseline{150}{他们为何造起这匹巨马？谁出的主意？}
\verseline{}{所求何事？是祭仪，还是战争的器械？’}
\verseline{}{那人深谙佩拉斯吉人的诡计，}
\verseline{}{把挣脱锁链的双手举向群星：}
\verseline{}{‘永恒的火焰啊，我以你们不可侵犯的}
\verseline{155}{神威为证，’他说，‘也以我逃离的祭坛和可怖的刀剑，}
\verseline{}{以及我作为祭牲戴过的众神圣带为证。\footnote{vittae 是祭牲或祭司佩戴的圣带，与普通装饰头带有别。西农以自己逃离的祭坛、刀剑和圣带发誓，借这些祭仪物件把自己说成原本将被献祭的牺牲者。}}
\verseline{}{我有权解开希腊人神圣的誓约，}
\verseline{}{有权憎恨他们，把他们遮掩的一切}
\verseline{}{暴露在天光下。祖国的任何律法都不再约束我。}
\verseline{160}{只要你守住所许，只要特洛伊得救以后}
\verseline{}{仍守信于我——若我所言是真，所报重大。}
\verseline{}{达那俄斯人全部的希望与开战的信心}
\verseline{}{始终仰赖帕拉斯的援助。可是自从那渎神的}
\verseline{}{提丢斯之子和罪恶的发明者乌利塞斯，}
\verseline{165}{胆敢从圣庙中拔走那命定的}
\verseline{}{帕拉斯神像，杀死城堡顶端的守卫，\footnote{“提丢斯之子”是狄俄墨得斯。他与乌利塞斯盗走帕拉斯神像（Palladium）的故事不见于《伊利亚特》正文，却是后来的特洛伊传说核心；西农据此谎称女神已离弃希腊人。}}
\verseline{}{抢走神圣的形象，又用染血的双手}
\verseline{}{触碰女神处女般纯洁的圣带；}
\verseline{}{从那时起，达那俄斯人的希望便开始退落}
\verseline{170}{力量摧折，女神的心也已背转。}
\verseline{}{特里同尼亚以无可置疑的异兆显明此事。}
\verseline{}{神像刚在营中安放，她高昂的双眼里}
\verseline{}{便燃起闪烁的火焰，咸涩的汗水流遍四肢；}
\verseline{}{三次，她从地上——说来可惊——}
\verseline{175}{猛然跃起，手持圆盾，长矛颤动。}
\verseline{}{卡尔卡斯立刻宣称必须冒险逃向海上，}
\verseline{}{阿尔戈斯人的兵器再也不能摧毁帕加马，}
\verseline{}{除非他们回到阿尔戈斯重新求取征兆，把那神威}
\verseline{}{带回来——如今他们已用弯曲的船载它渡海而去。}
\verseline{180}{所以他们乘风回到祖先的迈锡尼，}
\verseline{}{正在备好武器，邀诸神同行，等重新越海，}
\verseline{}{便会骤然归来。卡尔卡斯如此解释征兆。}
\verseline{}{他们奉警示造起这座木像，代替帕拉斯神像，}
\verseline{}{以赎偿那被冒犯的神威和悲惨的渎神之罪。}
\verseline{185}{但卡尔卡斯命令他们把它造得无比庞大}
\verseline{}{以交织的橡木高高举向天空，}
\verseline{}{使它不能被城门容纳，也不能拖进城墙，}
\verseline{}{不让它依照古老信仰庇护你们的人民。}
\verseline{}{因为倘若你们的手亵渎弥涅耳瓦的礼物，}
\verseline{190}{巨大的毁灭便会降临——愿诸神先把这个凶兆}
\verseline{}{转到他自己身上！——落在普里阿摩斯的王权和佛里吉亚人身上；}
\verseline{}{但若它由你们的手升入你们的城市，}
\verseline{}{亚细亚便会主动发动大战，直抵珀罗普斯的城墙，}
\verseline{}{这样的命数将等待我们的后代。’}
\verseline{195}{西农凭这样的诡计和伪誓手腕}
\verseline{}{使人信服，我们被骗局和伪造的眼泪俘获——}
\verseline{}{提丢斯之子、拉里萨的阿喀琉斯未能征服我们，}
\verseline{}{十年和一千艘战船也未能征服，如今我们竟就此被俘。}
\verseline{}{这时又一件更大、远更可怖的怪事\footnote{“更大、远更可怖的怪事”指后文自海中突现的双蛇异象（2.203—227）。此前西农的巧辩（2.67—194）已动摇特洛伊人，而这一骇人神兆的骤降彻底瓦解了众人的戒心。拉奥孔作为涅普顿祭司正宰杀公牛，却随即与二子一同沦为巨蛇的牺牲品。}}
\verseline{200}{降临到我们这些不幸者眼前，扰乱毫无防备的心。}
\verseline{}{拉奥孔由抽签选作涅普顿的祭司，}
\verseline{}{正在庄严的祭坛前宰杀一头巨牛。}
\verseline{}{看，两条巨蛇从忒涅多斯而来，越过平静的深海——}
\verseline{}{说到这里我仍会战栗——以庞大的环身}
\verseline{205}{压住海面，并排向岸边游来。}
\verseline{}{它们的胸膛在浪间高昂，血红冠脊}
\verseline{}{升出波涛，其余身躯在后方掠过海面，}
\verseline{}{巨大的脊背蜿蜒盘曲。}
\verseline{}{泡沫翻涌的盐海轰然作响。它们已经登岸，}
\verseline{210}{双眼浸满鲜血和火焰，灼灼燃烧，}
\verseline{}{颤动的舌头舔过嘶鸣的口。}
\verseline{}{我们一见便面无人色，四散奔逃。它们却不偏不倚，}
\verseline{}{直取拉奥孔。两条蛇先各自缠住}
\verseline{}{他两个儿子幼小的身体，}
\verseline{215}{盘结起来，啃噬可怜孩子的肢体。}
\verseline{}{随后拉奥孔持械奔来救援，}
\verseline{}{它们又把他攫住，用巨大的蛇环捆缚，}
\verseline{}{两度绕住腰身，两度以带鳞的脊背}
\verseline{}{缠过颈项，头颅和高昂的脖颈仍凌驾其上。}
\verseline{220}{他双手拼命撕扯那些死结，}
\verseline{}{圣带浸透污血和黑色毒液，}
\verseline{}{同时向群星发出骇人的呼号：}
\verseline{}{恰如一头受伤的公牛吼叫着逃离祭坛，}
\verseline{}{甩脱落在颈上、没有砍准的斧头。}
\verseline{225}{两条巨蛇却滑向神庙最高处}
\verseline{}{逃进残酷的特里同尼亚的城堡，}
\verseline{}{藏在女神脚下，她圆盾的遮蔽之中。}
\verseline{}{于是新的恐惧潜入所有人战栗的胸膛}
\verseline{}{他们说拉奥孔已为罪行付出应得的代价，\footnote{两条海蛇从忒涅多斯方向游来，先杀拉奥孔二子，再缠死祭司本人。特洛伊人因而把拉奥孔之死理解为他投矛冒犯木马的惩罚。}}
\verseline{230}{因为他以矛尖伤害神圣的木身，}
\verseline{}{把那不敬的长矛掷进木马背上。}
\verseline{}{众人齐声高喊，要把木像拖到它的居所}
\verseline{}{向女神的神威祈祷。}
\verseline{}{我们拆开城墙，敞开城市的防御。}
\verseline{235}{所有人都投入劳作，在木马脚下}
\verseline{}{垫入滚动的轮架，把麻绳}
\verseline{}{套紧它的颈项。致命的器械爬上城墙，}
\verseline{}{腹中怀着兵器。男孩和未婚少女环绕四周，}
\verseline{}{唱着圣歌，欣喜地用手触碰绳索。}
\verseline{240}{它越过门槛，带着威胁滑进城心。}
\verseline{}{哦祖国，哦众神的家园伊利乌姆，哦达耳达诺斯人}
\verseline{}{以战争闻名的城墙！它四次停在城门}
\verseline{}{门槛上，腹中的兵器也四次发出声响。}
\verseline{}{我们却仍旧坚持，浑然忘却一切，狂热蒙住双眼，}
\verseline{245}{终于把这不祥的怪物安置在神圣城堡。}
\verseline{}{那时卡珊德拉也曾开口揭示未来的命数}
\verseline{}{但依神明的命令，特乌克罗斯人从不相信她。\footnote{卡珊德拉受阿波罗赐予预言能力，却因拒绝神明而被诅咒永不获信。她在胜利庆典中开口，说出的真话仍无人相信。}}
\verseline{}{我们这些不幸者——那一天本是最后一天——}
\verseline{}{却还用节庆的枝叶装饰全城神殿。}
\verseline{250}{与此同时天空转动，黑夜从海上涌来，}
\verseline{}{以巨大的阴影笼罩大地与天穹，}
\verseline{}{也笼罩密耳弥冬人的诡计。特乌克罗斯人散卧城中，}
\verseline{}{尽皆无声，睡眠拥住疲惫的四肢。}
\verseline{}{阿尔戈斯的方阵已乘着整备好的船只}
\verseline{255}{从忒涅多斯出发，借静月友善的沉默}
\verseline{}{驶向熟悉的海岸。王船举起火焰信号，}
\verseline{}{西农仗着众神不义的命数护持}
\verseline{}{暗中放出腹中的达那俄斯人，并将松木门闩}
\verseline{}{松开。木马向空气敞开，把他们送出；}
\verseline{260}{领袖们欣然从空洞的木身中现身：}
\verseline{}{忒散德罗斯、斯忒涅卢斯和可怖的乌利塞斯，}
\verseline{}{顺着垂下的绳索滑落，还有阿卡玛斯、托阿斯，}
\verseline{}{珀琉斯之孙涅俄普托勒摩斯、最先下来的玛卡翁，}
\verseline{}{墨涅拉俄斯，以及诡计的制造者厄佩俄斯本人。}
\verseline{265}{他们侵入被睡眠和美酒埋葬的城市；}
\verseline{}{杀死守卫，洞开城门迎入}
\verseline{}{全部同伴，与知晓阴谋的队伍会合。}
\verseline{}{那正是劳苦已久的凡人刚刚\footnote{本卷写希腊舰队从忒涅多斯驶回，王船举火，西农放出木马伏兵（2.254—259）。阿波罗多洛斯《书库·摘要》5.15、5.19另有西农举火示意的传统。海伦绕木马、模仿诸将妻子声音，见《奥德赛》4.271—289；本诗6.518—519由得伊福玻斯说她举火召来希腊人。}}
\verseline{}{沉入初睡之时——众神最可喜的赠礼缓缓爬行。}
\verseline{270}{梦中，看啊，悲痛至极的赫克托耳}
\verseline{}{似乎来到我眼前，泪流不止，}
\verseline{}{仍像从前被双马战车拖曳，浑身发黑}
\verseline{}{带血的尘土沾满全身，肿胀的双脚被缰索穿透。}
\verseline{}{唉，他成了何等模样！与那个赫克托耳相比何其不同——}
\verseline{275}{那个曾披着阿喀琉斯铠甲归来，}
\verseline{}{或把佛里吉亚的火焰掷向达那俄斯战船的人。}
\verseline{}{如今胡须蓬乱，发丝凝着鲜血}
\verseline{}{带着绕祖国城墙时承受的}
\verseline{}{无数伤口。我自己仿佛也在流泪}
\verseline{280}{主动向他开口，吐出悲伤的话：}
\verseline{}{‘哦达耳达诺斯的光，哦特乌克罗斯人最可靠的希望，\footnote{赫克托耳以被阿喀琉斯拖尸后的伤残形貌入梦，昔日英雄盛装已然消失。他把城的圣物与家神交给埃涅阿斯，标志守城使命已经结束；接下来埃涅阿斯要把圣物、家神和幸存者带出城。马克罗比乌斯（6.2.18）还保存恩尼乌斯悲剧《亚历山大》一句 O lux Troiae ... Hector，与这里 O lux Dardaniae ... 的开头紧密对应。}}
\verseline{}{什么使你耽搁这样久？久盼的赫克托耳，}
\verseline{}{你从哪片海岸归来？同胞已经死去如此多人}
\verseline{}{人民与城市已经历尽种种劳苦，}
\verseline{285}{疲惫的我们终于看见了你！什么不堪的缘故}
\verseline{}{污损你安详的面容？我为什么看见这些伤口？’}
\verseline{}{他不作回答，也不理会我徒劳的询问，}
\verseline{}{只从胸膛最深处拖出沉重的叹息：}
\verseline{}{‘唉，逃吧，女神之子，快从火焰中救出自己。}
\verseline{290}{敌人占据了城墙，特洛伊正从高高的巅峰坠落。}
\verseline{}{你为祖国、为普里阿摩斯付出的已经够多。若有一只右手}
\verseline{}{能够守住帕加马，它早已被这只手守住。}
\verseline{}{特洛伊把圣物和自己的家神托付给你；}
\verseline{}{带上这些命运的同伴，为他们寻找宏伟的城墙——}
\verseline{295}{越过大海漂泊之后，你终将把它建起。’}
\verseline{}{说罢，他亲手从内殿深处捧出圣带、}
\verseline{}{强大的维斯塔与不灭的火。\footnote{维斯塔是炉火与城邦延续的女神；不灭之火与家神、圣带一同被带出内殿。埃涅阿斯带走的不只是幸存者，也包括特洛伊的公共祭祀与家族宗教传统。}}
\verseline{}{与此同时城中各处哀声四起，}
\verseline{}{声音越来越清楚。尽管我父亲}
\verseline{300}{安喀塞斯的家远在僻处，被树木遮蔽，}
\verseline{}{兵器的恐怖仍一阵阵逼近。}
\verseline{}{我从睡梦中惊起，朝屋顶最高处}
\verseline{}{攀爬而上，侧耳站定：}
\verseline{}{恰如狂暴的南风把火焰吹进麦田，}
\verseline{305}{或山河中迅猛的洪流，}
\verseline{}{扫平田野、茂盛的庄稼和牛耕的成果，}
\verseline{}{把森林倒卷而去，牧人茫然站在}
\verseline{}{高岩之上，听见轰声便惊得失神。}
\verseline{}{这时真相终于确凿，达那俄斯的诡计}
\verseline{310}{全然显露。得伊福玻斯宽大的屋宇\footnote{得伊福玻斯是普里阿摩斯之子，帕里斯死后海伦嫁给他；其遭难详见第六卷。乌卡勒贡是特洛伊长老。第312行以他的人名代指住宅。}}
\verseline{}{已被武尔坎努斯的烈火压倒轰然崩塌，近旁，}
\verseline{}{乌卡勒贡的宅邸也在燃烧，西格乌姆的广阔海峡映着火光。}
\verseline{}{众人的呼号和号角的鸣响一同升起。}
\verseline{}{我发狂般抓起武器。拿起武器已谈不上理智，}
\verseline{315}{但胸中燃烧，只想召集一队人}
\verseline{}{同伙伴冲向城堡。狂热和愤怒把我}
\verseline{}{推向前去，我只觉得持械而死也很壮美。}
\verseline{}{看，潘图斯刚从阿凯亚人的兵器中逃出——}
\verseline{}{俄特律斯之子，城堡和福玻斯的祭司，}
\verseline{320}{圣物、战败的神明和年幼的孙儿都在他手中，}
\verseline{}{他亲自携着这一切，发狂似地奔到我门前。}
\verseline{}{‘大势如何，潘图斯？我们据守哪座城堡？’}
\verseline{}{我刚说完，他便在呻吟中这样回答：}
\verseline{}{‘最后的日子已经来到，达耳达诺斯无可逃避的}
\verseline{325}{时刻已经来到。我们曾是特洛伊人，伊利乌姆曾经存在；}
\verseline{}{特乌克罗斯人曾有宏大的荣耀。残酷的朱庇特已把一切}
\verseline{}{转给阿尔戈斯，达那俄斯人在燃烧的城中称王。}
\verseline{}{高耸的木马立在城中央，正倾吐}
\verseline{}{全副武装的士兵。胜利的西农傲慢地四处纵火。}
\verseline{330}{另一些人从洞开的双扇城门涌入，}
\verseline{}{人数多得像伟大的迈锡尼曾派来的全部军队；}
\verseline{}{另一些人持械堵住狭窄的街道。}
\verseline{}{铁刃出鞘，剑尖寒光闪动}
\verseline{}{已经备好杀戮。城门最前的守卫才刚刚接战，}
\verseline{335}{在混乱不明的搏斗中抵抗。’}
\verseline{}{听完俄特律斯之子这番话，又受众神之力驱使，}
\verseline{}{我奔向火焰与兵器，去往悲厉的复仇女神}
\verseline{}{与喧声、直冲天穹的呼喊召唤我的地方。}
\verseline{}{瑞剖斯也来加入，武艺最强的厄皮图斯}
\verseline{340}{借月光出现，许帕尼斯和底玛斯也一同前来，}
\verseline{}{聚到我身旁，还有年轻的科罗厄布斯，}
\verseline{}{弥格冬之子。那些日子他恰好来到特洛伊，}
\verseline{}{因对卡珊德拉的痴狂爱情而燃烧，}
\verseline{}{以未来女婿身份援助普里阿摩斯和佛里吉亚人，}
\verseline{345}{不幸的人竟没有听从发狂的未婚妻}
\verseline{}{发出的警告。}
\verseline{}{我见他们聚在一处，仍敢投入战斗，}
\verseline{}{便接着说道：‘年轻人啊，你们这徒然勇敢的心！}
\verseline{}{倘若你们已决意随我冒险，}
\verseline{350}{一直走到绝境，你们也看见了大势如何：}
\verseline{}{那些曾支撑这座王国的神明已经全部离去，}
\verseline{}{抛下神殿深处和祭坛。你们正赶来}
\verseline{}{援救一座燃烧的城市。让我们冲进兵刃之间赴死吧。}
\verseline{}{战败者唯一的生路，就是不指望任何生路。’}
\verseline{355}{这番话又把狂热注入青年们的心。于是如同}
\verseline{}{黑雾中出巢掠食的群狼，腹中饥饿}
\verseline{}{凶暴地逼迫它们，而留在巢穴的幼崽}
\verseline{}{正喉咙干渴地等着，我们也这样穿过投枪与敌人}
\verseline{}{走向必死的结局，沿城市中央道路}
\verseline{360}{前行，黑夜以空洞的阴影环绕我们。}
\verseline{}{谁能用言语展开那一夜的灾祸，谁能说尽死者，}
\verseline{}{又有谁的泪水足以与那些苦难相称？}
\verseline{}{古老的城市倾覆了，它曾统治许多年；}
\verseline{}{无数失去生命的躯体到处横陈，散落在街道}
\verseline{365}{和屋内，也倒在众神神圣的门槛上。}
\verseline{}{也并非只有特乌克罗斯人以鲜血偿还；}
\verseline{}{有时勇气也重返战败者的胸膛}
\verseline{}{胜利的达那俄斯人也随之倒下。处处是残酷的哀恸}
\verseline{}{处处是恐惧，死亡以无数形貌出现。}
\verseline{370}{安德罗革俄斯最先迎面出现，身后跟着大群}
\verseline{}{达那俄斯人。他毫不知情，以为我们是友军}
\verseline{}{竟主动用亲切的话招呼：}
\verseline{}{‘快些，弟兄们！怎么耽搁到这样晚？}
\verseline{}{别人都在洗劫燃烧的帕加马，搬走战利品，}
\verseline{375}{你们现在才从高大的战船下来。’}
\verseline{}{他说完，因给出的回答并不够可信，}
\verseline{}{便立刻察觉自己已落入敌人中央。}
\verseline{}{他惊住，话音和脚步一同收回：}
\verseline{}{恰如有人毫无防备地穿过粗硬的荆棘，}
\verseline{380}{一脚踩住地上的蛇，惊惧中骤然退开，}
\verseline{}{蛇正昂起怒气，鼓胀青蓝色的颈项。}
\verseline{}{安德罗革俄斯望见我们，也这样战栗着退去。}
\verseline{}{我们扑上去，用密集的兵器包围他们，}
\verseline{}{他们不熟悉地势，又被恐惧攫住，}
\verseline{385}{我们便四下将他们击倒。命运眷顾我们的第一次交锋。}
\verseline{}{科罗厄布斯因得手而振奋，欢喜地说：}
\verseline{}{‘同伴们，命运最先在哪里向我们指出一条}
\verseline{}{得救的道路，又在哪里显出有利的一面，我们就跟到哪里。}
\verseline{}{换掉盾牌，把达那俄斯人的徽记}
\verseline{390}{披到自己身上。面对敌人，谁会追究是诡计还是勇气？}
\verseline{}{他们自己会把武器给我们。’说罢，他穿戴起}
\verseline{}{安德罗革俄斯的羽冠头盔和饰有华丽徽记的盾牌，}
\verseline{}{又把阿尔戈斯长剑佩在腰间。}
\verseline{}{瑞剖斯、底玛斯和所有青年也都如此}
\verseline{395}{人人欢喜地用新获的战利品武装自己。}
\verseline{}{我们混进达那俄斯人中，借一份不属于我们的神佑前行，}
\verseline{}{在盲目的黑夜里多次交锋}
\verseline{}{与敌人厮杀，把许多达那俄斯人送入奥尔库斯。}
\verseline{}{有些人四散奔向战船，以最快的脚步}
\verseline{400}{寻找可靠的海岸。另一些人因可耻的恐惧，}
\verseline{}{又爬回木马，藏进熟悉的腹腔。}
\verseline{}{唉，若众神不愿，凡人便无权依恃任何事！\footnote{nihil invitis fas quemquam fidere divis 说：诸神不许时，凡人便无神律许可去依靠任何事。埃涅阿斯由此承认，冒用希腊兵甲的计策缺少神意支持。}}
\verseline{}{看啊，普里阿摩斯的女儿卡珊德拉披散头发，}
\verseline{}{正被人从弥涅耳瓦的神庙、从内殿拖出，}
\verseline{405}{徒然把燃烧的双眼举向天空——}
\verseline{}{只能举眼，因为绳索缚住了她柔嫩的手掌。\footnote{卡珊德拉被从弥涅耳瓦内殿拖出，呼应小埃阿斯亵渎她的传统。科罗厄布斯是她的求婚者；救她时，他与同伴因所穿希腊兵甲遭己方误击。}}
\verseline{}{科罗厄布斯狂乱的心无法承受这景象，}
\verseline{}{纵身投入敌阵，注定要死在那里。}
\verseline{}{我们全都跟上，用密集的兵器冲入敌阵。}
\verseline{410}{这时神庙高顶上的同胞先向我们}
\verseline{}{投下武器。最惨烈的杀戮由此发生，}
\verseline{}{只因盔甲的外貌，只因希腊羽冠造成的误认。}
\verseline{}{达那俄斯人又为少女被夺而怒吼}
\verseline{}{从四面聚来进攻：最凶猛的埃阿斯，}
\verseline{415}{阿特柔斯二子，以及多洛佩斯人的全军。}
\verseline{}{恰如风暴一旦破裂，迎面的狂风}
\verseline{}{彼此冲撞：西风、南风，还有驾着东方骏马}
\verseline{}{欢腾而来的东风。森林尖啸，涅柔斯满身泡沫}
\verseline{}{挥动三叉戟，从海底搅起大水。}
\verseline{420}{那些先前被我们借黑夜的阴影}
\verseline{}{设伏击溃、追遍全城的人也重新出现，}
\verseline{}{他们首先认出虚假的盾牌和兵器，}
\verseline{}{又从不相谐的口音认出我们的身份。}
\verseline{}{我们立刻被人数压倒。科罗厄布斯第一个}
\verseline{425}{倒在佩涅琉斯右手之下、善战女神的祭坛旁，}
\verseline{}{瑞剖斯也倒下了——在特乌克罗斯人中，\footnote{瑞剖斯被称为特乌克罗斯人中“最公正、最守正”的一人，却仍在当夜倒下；诗人随即说诸神另有判断。他的正直没有使他免于死亡。}}
\verseline{}{他最为公正，也最恪守正道：}
\verseline{}{众神却另有裁断。许帕尼斯和底玛斯}
\verseline{}{被自己同伴的兵器刺穿。潘图斯啊，}
\verseline{430}{你倒下时，再深的虔敬与阿波罗的圣带也未能护你。}
\verseline{}{伊利乌姆的灰烬啊，我族最后的火焰，}
\verseline{}{我以你们作证：你们覆亡之时我没有躲避}
\verseline{}{达那俄斯人的兵器，也没有避开任何交锋；}
\verseline{}{倘若命运要我倒下，我的手也已挣得死亡。我们随即被冲散，}
\verseline{435}{同我一起的是伊菲图斯和珀利阿斯。伊菲图斯}
\verseline{}{已为年岁所累，珀利阿斯又因乌利塞斯所伤而行动迟缓；}
\verseline{}{一阵呼喊随即把我们召向普里阿摩斯的宫殿。}
\verseline{}{在那里我们才看见一场巨大的战斗，仿佛别处}
\verseline{}{根本没有战争，整座城也无人死去。}
\verseline{440}{我们看见不驯的战事，看见达那俄斯人冲向屋宇；}
\verseline{}{大门已被推进的龟甲阵围困。}
\verseline{}{云梯贴住墙壁，他们就在门柱下方}
\verseline{}{攀着梯级上升，左手以盾牌迎向投枪}
\verseline{}{护住身体，右手抓住屋檐。}
\verseline{445}{达耳达诺斯人则从上方拆下塔楼和整片}
\verseline{}{屋顶。他们看见末日已至，准备用这些}
\verseline{}{兵器在死亡尽头守卫自己；}
\verseline{}{推下镀金的横梁——古老先祖的华贵装饰。}
\verseline{}{另一些人拔剑围住下方门户}
\verseline{450}{结成密集的队列把守。}
\verseline{}{我们的心重新振起，要援救国王的屋宇，}
\verseline{}{替守卫者减轻压力，给战败者增添力量。}
\verseline{}{那里有一道门，一处隐蔽而畅通的入口}
\verseline{}{连接普里阿摩斯各座宫室，后方留着一扇门；}
\verseline{455}{王国尚存时，不幸的安德罗玛刻常从这里}
\verseline{}{独自前往公婆的家中，}
\verseline{}{又牵着男孩阿斯提阿那克斯去见祖父。}
\verseline{}{我爬上最高屋脊。可怜的特乌克罗斯人}
\verseline{}{正在那里向下投出徒劳的兵器。}
\verseline{460}{一座塔楼立在陡峭边缘，从屋顶最高处}
\verseline{}{高耸入星空。由那里，整个特洛伊、}
\verseline{}{达那俄斯人的船只与阿凯亚营地，往常都尽收眼底。}
\verseline{}{我们用铁器从四周下手，在宫殿顶层的楼板}
\verseline{}{露出松动接缝之处，将塔楼从高处的基座撬开}
\verseline{465}{合力推下。塔楼突然坠落，带着整片崩塌，}
\verseline{}{轰然砸进达那俄斯人的队列。}
\verseline{}{但另一些人又冲上来，石块和各种}
\verseline{}{兵器都不曾停歇。}
\verseline{}{就在前庭、第一道门槛前，皮洛士\footnote{皮洛士即涅俄普托勒摩斯，阿喀琉斯之子；此处随拉丁原文换称。他闯入普里阿摩斯宫殿时多次被写得像父亲阿喀琉斯，最后却在祭坛旁杀死老人普里阿摩斯。}}
\verseline{470}{身披青铜，在兵光中跳跃逼近；}
\verseline{}{如同一条饱食毒草、才爬到阳光下的蛇，}
\verseline{}{寒冷的冬季曾让它肿胀着藏在地下，}
\verseline{}{如今脱去旧皮，青春焕发，}
\verseline{}{昂起胸膛，盘卷光滑的脊背}
\verseline{475}{迎向太阳，三叉的舌尖在口中闪动。}
\verseline{}{同他一起的有高大的佩里法斯，还有阿喀琉斯的御者，}
\verseline{}{即武器侍从奥托墨冬，以及全部斯库罗斯青年，}
\verseline{}{他们逼近宫殿，把火焰投向屋顶。}
\verseline{}{他冲在最前，抓起双刃斧，}
\verseline{480}{劈破坚固的门户，从门轴处猛扯包铜的门柱，}
\verseline{}{又砍断横梁，凿穿坚实的橡木}
\verseline{}{在门上豁开一道宽阔的巨口。}
\verseline{}{屋内显露出来，长长的中庭豁然敞开；}
\verseline{}{普里阿摩斯和古代诸王的内室也暴露无遗，}
\verseline{485}{他们看见武装的人守在最前一道门槛。}
\verseline{}{宫殿深处一片混乱，呻吟四起，}
\verseline{}{空荡的屋宇从最深处回响女人捶胸的哀号，}
\verseline{}{呼声撞向金色的群星。}
\verseline{}{惊恐的母亲们在巨大的宫室间游荡；}
\verseline{490}{抱住门柱，把亲吻印在上面。}
\verseline{}{皮洛士凭父亲传下的力量猛逼，门闩和守卫}
\verseline{}{都挡不住他。大门在接连撞击下摇晃，}
\verseline{}{门柱挣脱铰链，向前倒去。}
\verseline{}{以强力开路；达那俄斯人破门涌入}
\verseline{495}{杀倒最前列的人，士兵顷刻填满四周。}
\verseline{}{决堤的河流也不曾这样凶猛。泡沫翻滚的水}
\verseline{}{冲破堤岸，用漩涡压倒迎面的障碍，}
\verseline{}{堆叠着奔向田野，席卷所有平原}
\verseline{}{连牛群与棚圈一并拖走。我亲眼看见}
\verseline{500}{杀戮中发狂的涅俄普托勒摩斯，也看见门槛上的阿特柔斯二子；}
\verseline{}{看见赫库芭与百名儿媳，也看见普里阿摩斯在祭坛之间}
\verseline{}{用鲜血玷污自己亲手祝圣的火焰。}
\verseline{}{那五十间婚房——后代如此众多的希望——}
\verseline{}{那些以异邦黄金和战利品自豪的门柱}
\verseline{505}{尽皆倒下。火焰止步之处也被达那俄斯人占据。}
\verseline{}{也许你还想问，普里阿摩斯遭逢了怎样的命运。\footnote{赫库芭与女儿们聚在宫中祭坛旁；普里阿摩斯披上久不用的兵甲后仍被妻子拉回祭坛。祭坛本应提供庇护，却最终成为皮洛士杀死普里阿摩斯的地点。第521行的“守卫”亦可指普里阿摩斯身上的兵甲（Fairclough—Brown，\textit{Virgil’s Aeneid: Books I–VI}，1908，评2.521；\href{https://dcc.dickinson.edu/vergil-aeneid/vergil-aeneid-ii-506-525}{经 DCC 汇录}）。更早的拉丁悲剧记忆也压在这一场景上：恩尼乌斯《安德罗玛刻》残篇让说话者“亲眼看见”特洛伊焚毁、普里阿摩斯被杀、朱庇特祭坛被鲜血玷污；维吉尔在 2.499--502 以 vidi、Priamum、per aras／sanguine foedantem 几乎逐项重组这组画面。其中“鲜血玷污”的措辞又叠着卢克莱修《物性论》1.84--85 伊菲革涅亚祭坛场景的语言。}}
\verseline{}{当他看见城市陷落的灾难，看见宫门}
\verseline{}{被拔毁，敌人已经进入内殿中央，}
\verseline{}{老人徒然把久已不用的武器披上}
\verseline{510}{因年老而颤抖的双肩，又佩起无用的铁剑}
\verseline{}{冲向密集的敌人，准备赴死。}
\verseline{}{宫室中央，在露天的苍穹之下}
\verseline{}{立着一座巨大的祭坛，旁边有棵极古老的月桂树，}
\verseline{}{俯临祭坛，以树荫环抱家神。}
\verseline{515}{赫库芭和女儿们徒然围在祭坛四周，}
\verseline{}{如黑色风暴中急坠的鸽群，}
\verseline{}{她们挤坐在一起，抱紧众神的神像。}
\verseline{}{赫库芭看见普里阿摩斯竟穿上青年时代的武器}
\verseline{}{便说：‘我最不幸的丈夫，什么可怕的念头}
\verseline{520}{驱使你佩起这些兵器？你要奔向哪里？}
\verseline{}{此刻不需要这样的援助，也不需要这样的守卫；}
\verseline{}{即使我的赫克托耳亲自在这里也无济于事。}
\verseline{}{快退到这里来吧。这座祭坛将庇护我们所有人，}
\verseline{}{否则你就同我们一起死。’她说着，把老人}
\verseline{525}{拉到身旁，安置在神圣的座位上。}
\verseline{}{看，波利特斯逃过皮洛士的屠杀，}
\verseline{}{这位普里阿摩斯之子穿过兵器、穿过敌人，}
\verseline{}{带着伤沿长廊奔逃，越过空荡的中庭，}
\verseline{}{皮洛士怒火中烧，挺着敌意的长矛，}
\verseline{530}{紧追不舍，眼看就要抓住他，枪尖已抵近身后。}
\verseline{}{当他终于逃到父母眼前，}
\verseline{}{便倒在那里，生命随着涌出的鲜血流尽。}
\verseline{}{普里阿摩斯虽已被死亡围在中央，}
\verseline{}{仍不肯忍住，也不吝惜声音和愤怒：}
\verseline{535}{‘愿诸神因你的罪行和这般胆大妄为，’他喊道，}
\verseline{}{‘倘若天上尚有敬义，肯看顾这样的事，}
\verseline{}{便向你偿付相称的报答，给你应得的奖赏。}
\verseline{}{你迫使我亲眼看见儿子被杀}
\verseline{}{以他的死亡玷污了父亲的面容。}
\verseline{540}{可那位你虚称为父的阿喀琉斯}
\verseline{}{对敌人普里阿摩斯也不曾如此。他敬畏}
\verseline{}{祈求者的权利和信义，把赫克托耳苍白的遗体}
\verseline{}{还给坟墓，又送我返回自己的王国。’}
\verseline{}{老人说罢，掷出一支软弱无力的长矛，}
\verseline{545}{它撞上青铜发出低哑声响，立刻弹开，}
\verseline{}{徒然挂在盾牌最高的凸脐上。}
\verseline{}{皮洛士答道：‘那么你就去做信使}
\verseline{}{把这些话带给我父亲、珀琉斯之子。记得向他讲述}
\verseline{}{我的残酷行径，讲述涅俄普托勒摩斯怎样辱没先人。}
\verseline{550}{现在死吧。’说着他把战栗的老人拖到祭坛前}
\verseline{}{老人在儿子流出的血泊中滑跌，}
\verseline{}{左手揪住头发，右手高举长剑，}
\verseline{}{寒光刺进肋侧，剑刃直没至柄。}
\verseline{}{这便是普里阿摩斯命运的终点。这样的结局}
\verseline{555}{带走了他——他亲眼看见特洛伊焚烧、帕加马倾覆，}
\verseline{}{他曾统治亚细亚，为众多民族和疆土骄傲。}
\verseline{}{海岸上横着一具巨大的无头躯干，}
\verseline{}{头颅从肩上撕去，只剩一具无名的身体。}
\verseline{}{这时一阵凶猛的恐惧第一次包围我。}
\verseline{560}{我惊呆了，亲爱父亲的形象浮上心头，}
\verseline{}{因为我看见一位与他同龄的国王，}
\verseline{}{从残酷的伤口中吐尽生命。我又想起孤立无援的克瑞乌萨，}
\verseline{}{想起被劫掠的家，想起幼小尤卢斯的命运。}
\verseline{}{我回头察看，环顾身边还剩多少人。}
\verseline{565}{所有人都疲惫地离开了。有些从屋顶，}
\verseline{}{把身体抛向地面，有些把虚弱的身躯交给火焰。}
\verseline{}{这时确实只剩我一人。我忽然在维斯塔的门槛旁，\footnote{第567—588行通常称为“海伦插曲”。Conte 2019 以斜体收入正文；该版校勘记指出早期主要抄本缺此段，主要由塞尔维乌斯评注传统保存。Conte 认为它可能未经诗人最后定稿，由诗人搁置或被早期编辑瓦里乌斯略去。}}
\verseline{}{看见廷达瑞俄斯之女，她藏在隐蔽处默不作声，}
\verseline{}{明亮的火光为我照明——我正四处游荡}
\verseline{570}{双眼不断扫过每一处地方。}
\verseline{}{她预先畏惧特乌克罗斯人因帕加马倾覆而仇视她}
\verseline{}{畏惧达那俄斯人的惩罚，也畏惧被遗弃丈夫的愤怒。}
\verseline{}{她是特洛伊和故国共同的复仇女神，}
\verseline{}{于是把自己藏起，可憎地坐在祭坛旁。}
\verseline{575}{我心中的火焰骤然燃起，愤怒袭来，要替}
\verseline{}{倾覆的祖国复仇，索取罪有应得的惩罚。}
\verseline{}{‘难道她竟要安然再见斯巴达和祖先的迈锡尼}
\verseline{}{以凯旋王后的身份离去，}
\verseline{}{再见丈夫、父亲的家宅和孩子，}
\verseline{580}{还有大群伊利乌姆妇女与佛里吉亚仆从随侍？}
\verseline{}{而普里阿摩斯死于铁剑，特洛伊焚于烈火？}
\verseline{}{达耳达诺斯的海岸又有多少次浸透鲜血？}
\verseline{}{绝不能如此。虽然惩罚一个女人}
\verseline{}{得不到值得传扬的名声，胜利也无荣耀，}
\verseline{585}{但消灭这场罪孽、向罪有应得者索偿，仍会使我受赞；}
\verseline{}{能以复仇的声名填满我的心，}
\verseline{}{又让同胞的灰烬得到满足，也会令我欣慰。’\footnote{Conte 2019 正文作 ultricis famae；该版校勘记另列 flammae 与 famam，并明言对正文读法仍有保留。ultricis 与 famae 在格、数、性上一致，意为“复仇的声名”；flammae 一读则为“复仇之火”。}}
\verseline{}{我喊出这些话，狂乱的心正把我带走。}
\verseline{}{这时慈爱的母亲亲自显现，从未如此清楚，}
\verseline{590}{她在黑夜中以纯净的光芒闪耀，}
\verseline{}{显露女神本相，带着她在天神面前惯常显出的}
\verseline{}{伟大身量。她用右手抓住我}
\verseline{}{拦下我，又用玫瑰色的唇说：}
\verseline{}{‘儿子，什么巨大的悲痛激起你这样无法约束的愤怒？}
\verseline{595}{你为何发狂？你对亲人的关切去了哪里？}
\verseline{}{你难道不先去看，年老疲惫的父亲}
\verseline{}{安喀塞斯被你留在何处，妻子克瑞乌萨，}
\verseline{}{和男孩阿斯卡尼乌斯是否还活着？希腊的阵列}
\verseline{}{正在他们四周游荡。若非我的守护抵挡，}
\verseline{600}{火焰早已带走他们，敌人的剑也早已吞噬他们。}
\verseline{}{不是那可憎的拉科尼亚女子的容貌，\footnote{维纳斯以诸神毁城之说制止埃涅阿斯向海伦复仇；她把特洛伊的覆亡归于诸神的无情。}}
\verseline{}{也不是受责的帕里斯；是众神，是众神的无情，}
\verseline{}{推倒了这片国势，使特洛伊从巅峰坠落。}
\verseline{}{看吧——我要从你眼前夺去那层遮蔽}
\verseline{605}{它正使凡人的目光迟钝，以潮湿的}
\verseline{}{阴翳笼住四周。不要害怕母亲的任何}
\verseline{}{命令，也不要拒绝服从我的告诫。}
\verseline{}{你看那里，巨构已经散裂，石块从石块上}
\verseline{}{崩脱，烟尘混合翻涌不息——}
\verseline{610}{是涅普顿以巨大的三叉戟撼动城墙}
\verseline{}{拔起根基，从最深的基座摧毁全城。}
\verseline{}{那里，最凶狠的朱诺首先占住斯开亚门，}
\verseline{}{又狂怒地从船边召来盟军，}
\verseline{}{披甲执兵。}
\verseline{615}{再回头看，特里同尼亚的帕拉斯已坐上最高城堡}
\verseline{}{在云光和凶恶的戈耳工中闪耀。}
\verseline{}{众神之父亲自给达那俄斯人添上勇气和得胜的力量}
\verseline{}{亲自激起诸神，反对达耳达诺斯的武器。}
\verseline{}{夺路逃走吧，儿子，给劳苦定下终点。}
\verseline{620}{我绝不会离开你，会把你安全送到父亲门前。’}
\verseline{}{她说完便藏入黑夜浓密的阴影。}
\verseline{}{可怖的神影显现——与特洛伊为敌的}
\verseline{}{伟大神祇。}
\verseline{}{这时我确实看见整个伊利乌姆沉入火焰}
\verseline{625}{看见涅普顿建造的特洛伊从根基倾覆；}
\verseline{}{恰如山顶一棵古老的白蜡树}
\verseline{}{已遭铁器砍伤，农人接连挥动双刃斧，}
\verseline{}{争相逼它倒下；它不断摇摇欲坠}
\verseline{}{枝冠震荡，树梢颤动，}
\verseline{630}{直到伤口逐渐将它征服，它发出最后一声}
\verseline{}{长吟，从山脊拔裂，拖着崩塌倒下。}
\verseline{}{我走下去，神明引路，使我穿过火焰和敌人}
\verseline{}{安然脱身。兵器让路，火焰退开。}
\verseline{}{当我终于回到祖屋古老的门槛}
\verseline{635}{父亲——我本想最先把他背上高山}
\verseline{}{也正是我最先奔来寻找的人——}
\verseline{}{却拒绝在特洛伊被摧毁后苟延生命}
\verseline{}{拒绝忍受流亡。他说：‘你们啊，血气仍在壮年}
\verseline{}{力量还凭自身坚实站立，}
\verseline{640}{你们赶快逃走。}
\verseline{}{倘若天神愿意让我继续生活，}
\verseline{}{他们原会为我保住这座家。一次毁灭已经太多，}
\verseline{}{我们看过城破，也从陷落的城市中活了下来。\footnote{安喀塞斯曾经历赫尔库勒斯在拉俄墨冬为王时攻破特洛伊。塞尔维乌斯（2.643） 将una excidia（一次毁灭）联系到这次旧事。}}
\verseline{}{就这样吧，把我当作已经躺下的尸体告别离去。}
\verseline{645}{我会亲手从敌人那里找到死亡。他会怜悯我}
\verseline{}{也会索取战利品。少一座坟墓也是容易承受的损失。}
\verseline{}{很久以来我便受诸神憎恶，徒然拖延}
\verseline{}{岁月。自从众神之父、人类之王}
\verseline{}{以雷霆之风吹中我，又以火焰触及我。’}
\verseline{650}{他反复说着这些，执意不移留在原处。}
\verseline{}{我们却泪流满面。妻子克瑞乌萨、}
\verseline{}{阿斯卡尼乌斯和全家一同恳求父亲不要}
\verseline{}{把一切同他一起拖入毁灭，不要向逼近的命运屈服。}
\verseline{}{他仍旧拒绝，固守原先的决定和住所。}
\verseline{655}{我这最不幸的人又奔向武器，渴望死亡：}
\verseline{}{因为还能有什么计策，什么命运留给我？}
\verseline{}{‘父亲，你竟以为我能留下你}
\verseline{}{独自迈步离开？这样巨大的不义怎会出自父亲之口？}
\verseline{}{倘若天上的神明决意不让这座伟大的城市留下一物，}
\verseline{660}{倘若你的心已经认定，要给垂死的特洛伊}
\verseline{}{再添上你和你的家人，那么死亡之门敞开着，}
\verseline{}{浑身带着普里阿摩斯鲜血的皮洛士很快就会来——}
\verseline{}{他当着父亲的面杀儿子，又在祭坛旁砍杀父亲。}
\verseline{}{慈爱的母亲啊，难道你把我从兵器和火焰中}
\verseline{665}{救出，只为让我在自己的内殿看见敌人}
\verseline{}{看见阿斯卡尼乌斯、我的父亲和身旁的克瑞乌萨}
\verseline{}{看他们一个倒在另一个人的血里，尽遭屠杀？}
\verseline{}{武器！战士们，拿武器来！末日召唤战败者。}
\verseline{}{把我送回达那俄斯人那里，让我再次冲进}
\verseline{670}{战斗。今天我们纵然全死，也绝不会无人替我们偿命。’}
\verseline{}{于是我又佩起铁剑，把左手}
\verseline{}{插进盾牌的带环，准备走出屋门。}
\verseline{}{看啊，妻子却在门槛前抱住我的双脚}
\verseline{}{紧贴不放，又把幼小的尤卢斯递向父亲：}
\verseline{675}{‘倘若你去送死，就把我们也带进一切险境。}
\verseline{}{倘若你身经百战，对手中武器仍有一点信心，}
\verseline{}{便先保护这个家。小尤卢斯，}
\verseline{}{你的父亲，还有我——曾被称作你妻子的人——你要留给谁？’}
\verseline{}{她这样高喊，呻吟充满整座屋宇，}
\verseline{680}{就在此时忽然出现一个说来惊异的征兆。}
\verseline{}{在悲伤父母的双手和面容之间}
\verseline{}{看啊，尤卢斯头顶的一簇微光似乎\footnote{尤卢斯头顶的火焰没有灼伤头发，随后又有雷声和流星确认。安喀塞斯先前拒绝逃离，直到连续征兆出现才改变决定；推动转折的是这组神示。}}
\verseline{}{倾泻光芒。柔软的火焰触之无害}
\verseline{}{舔过他的头发，在太阳穴周围燃烧。}
\verseline{685}{我们惊恐战栗，忙着拍打燃烧的头发}
\verseline{}{用泉水扑灭那神圣的火。}
\verseline{}{父亲安喀塞斯却欣喜地抬眼望向群星}
\verseline{}{向天空伸出双掌，高声说道：}
\verseline{}{‘全能的朱庇特，倘若任何祈祷能使你转意，}
\verseline{690}{请看我们——我只求这一件。若我们的虔敬当得起，}
\verseline{}{父亲，便赐下援助，确认这个征兆。’}
\verseline{}{老人刚说完，左方忽然雷声轰响，}
\verseline{}{一颗星从天空滑下，穿过幽暗，}
\verseline{}{拖着火炬般的长尾，在强光中疾行。}
\verseline{695}{我们看见它掠过屋顶最高的脊梁，}
\verseline{}{明亮地隐入伊达山的树林，}
\verseline{}{指明道路。随后长长的光痕}
\verseline{}{照亮轨迹，辽阔的四野弥漫硫黄烟气。}
\verseline{}{这时父亲终于折服，挺身站起，}
\verseline{700}{向众神开口，敬拜那颗神圣的星。}
\verseline{}{‘现在再无耽搁。我跟随。你们引向哪里，我便去哪里。}
\verseline{}{祖先的神明啊，保全这个家，保全我的孙儿。}
\verseline{}{这征兆属于你们，特洛伊也由你们的神意掌管。}
\verseline{}{我让步了，儿子，也不再拒绝与你同行。’}
\verseline{705}{他说完。城墙之内火焰的轰响已听得更加清楚，}
\verseline{}{大火把热浪卷得越来越近。}
\verseline{}{‘那么来吧，亲爱的父亲，伏到我的颈背上；}
\verseline{}{我会以双肩承受，这份劳苦不会使我沉重。}
\verseline{}{无论事情怎样落定，我们都面对同一场危险，}
\verseline{710}{也共有一条生路。让小尤卢斯}
\verseline{}{跟在我身边，妻子远远跟着我的脚步。}
\verseline{}{你们这些仆人，留心记住我要说的话。}
\verseline{}{城外有一座土丘和一座古老的神庙}
\verseline{}{那是荒废的刻瑞斯圣所，旁边有棵古柏，}
\verseline{715}{因祖先的敬奉已保存许多年。}
\verseline{}{我们会从不同方向来到这里会合。}
\verseline{}{父亲，你亲手带上圣物和祖先的家神；}
\verseline{}{我刚从这样巨大的战争和新近的杀戮中走出，}
\verseline{}{我不可触碰它们，必须先在活水中}
\verseline{720}{洗净自己。’}
\verseline{}{说罢我在低下的颈背和宽阔双肩上}
\verseline{}{铺好衣袍和黄褐色的狮皮，}
\verseline{}{俯身背起重负。小尤卢斯把小手\footnote{“小尤卢斯”与“阿斯卡尼乌斯”指同一人。安喀塞斯、孩子、家神和随后同行的妻子组成跨越三代的逃亡队列。}}
\verseline{}{扣在我的右手上，迈着赶不上父亲的小步跟随；}
\verseline{725}{妻子走在后面。我们穿过黑暗。\footnote{埃涅阿斯背父、携子、命妻随后而行，成为其 pietas 最著名的图像之一。克瑞乌萨随后的失散，使这幅建国祖先图像从一开始便带着不可补偿的缺口。}}
\verseline{}{不久前任何投来的兵器都不能动摇我}
\verseline{}{迎面结队而来的希腊人也不能使我畏惧，}
\verseline{}{如今每一阵风都令我惊恐，每一个声响都使我惊跳，}
\verseline{}{我的心高悬着，同样为身边的孩子和肩上的重负担忧。}
\verseline{730}{我已经接近城门，仿佛一路险境}
\verseline{}{都已逃过。突然密集的脚步声似乎}
\verseline{}{逼近耳边。父亲透过暗影向外望去，}
\verseline{}{大喊：‘儿子，快逃，儿子！他们逼近了。}
\verseline{}{我看见燃烧的盾牌和闪动的青铜！’}
\verseline{735}{这时不知哪一位与我为敌的神，在我惊惶时，}
\verseline{}{夺走了我混乱中的理智。}
\verseline{}{我奔入偏僻小路，离开熟悉的街道，}
\verseline{}{唉，可怜的我！妻子克瑞乌萨是遭命运夺去，}
\verseline{}{停留在后面，还是迷了路，或力气尽失、瘫坐下来——\footnote{Conte 2019 正文采用 lapsa；该版校勘记释作 languescens quia uiribus relicta，意为“因力气离身而衰弱”，可指克瑞乌萨力竭而坐下。这仍是埃涅阿斯关于她去向的猜测。}}
\verseline{740}{已无从知道。她再也没有回到我眼前。}
\verseline{}{直到我们来到古老刻瑞斯的土丘和圣所，\footnote{拉丁句法把克瑞乌萨的缺席延后到众人抵达刻瑞斯圣所、全部聚齐时才揭示；她的失踪甚至“瞒过”同伴、儿子和丈夫。延迟揭露加深了埃涅阿斯的惊骇。}}
\verseline{}{我才回头想到失去的她，才将心重新转向她。}
\verseline{}{众人终于在那里全部聚齐，唯独她一人缺席；}
\verseline{}{她的失踪竟瞒过了同伴、儿子和丈夫。}
\verseline{745}{我发狂时，没有哪个人、哪个神不曾遭我指责，}
\verseline{}{在倾覆的城市中，我还见过什么比这更残酷？}
\verseline{}{我把阿斯卡尼乌斯、父亲安喀塞斯和特乌克罗斯的家神\footnote{Teucros修饰Penatis，指特洛伊的家神。}}
\verseline{}{托付给同伴，让他们藏进曲折的山谷；}
\verseline{}{自己返回城市，再次佩起闪亮的武器。}
\verseline{750}{我决意重新承受全部风险，走遍}
\verseline{}{整个特洛伊，再一次把生命交给危险。}
\verseline{}{首先我回到城墙和那道幽暗的门槛，}
\verseline{}{就是我曾迈步离开的地方。我在夜中沿着}
\verseline{}{辨认出的足迹折返，以目光四处搜寻。}
\verseline{755}{恐怖处处攫住心灵，就连寂静本身也令人畏惧。}
\verseline{}{然后我回到家，看看她是否偶然}
\verseline{}{已经返回那里。达那俄斯人已经闯入，占满全屋。}
\verseline{}{吞噬一切的火焰立刻借风卷向最高屋顶}
\verseline{}{火舌凌越而上，热浪向天空狂怒。}
\verseline{760}{我继续前行，再去普里阿摩斯的宫殿和城堡。}
\verseline{}{朱诺圣所的柱廊已经空荡；}
\verseline{}{福尼克斯和可怖的乌利塞斯奉命守在那里，}
\verseline{}{看守战利品。特洛伊的珍宝从四方汇集到这里：}
\verseline{}{从燃烧的内殿抢出的物品、众神的供桌，}
\verseline{765}{纯金的调酒缸和被俘获的衣袍}
\verseline{}{堆成一处。男孩和惊恐的母亲，}
\verseline{}{排成长列站在四周。}
\verseline{}{我甚至敢于把呼声抛进阴影，}
\verseline{}{让呼喊填满街道，悲痛地呼唤克瑞乌萨，}
\verseline{770}{徒然一声又一声，一遍又一遍。}
\verseline{}{我寻找她，在城中屋宇间疯狂奔走不止，}
\verseline{}{不幸的幻影——克瑞乌萨本人的幽魂——}
\verseline{}{显现在眼前，比我熟悉的形体更加高大。}
\verseline{}{我惊住，头发直立而声音哽在喉中。}
\verseline{775}{于是她这样开口，用这些话减轻我的忧虑：\footnote{Conte 2019 将第775行置于方括号；该版校勘记记录 Ribbeck 的删疑，并指出同样的引语句式见第三卷第153行、第八卷第35行。克瑞乌萨幽魂预告埃涅阿斯将抵达赫斯珀里亚、获得王国与王族新娘；“众神伟大的母亲”通常指库柏勒。}}
\verseline{}{‘甜蜜的丈夫，这样放任疯狂的悲痛有什么用？}
\verseline{}{这一切并非没有众神的旨意。}
\verseline{}{神律不许你带克瑞乌萨作伴离开这里，}
\verseline{}{高踞奥林匹斯的统治者也不允许。}
\verseline{780}{漫长的流亡等待着你，辽阔的海面必须犁过；}
\verseline{}{你将来到赫斯珀里亚的土地，那里吕底亚的台伯河}
\verseline{}{以平缓的水流穿过人们富饶的田野。}
\verseline{}{在那里繁盛的命运、王国和一位王族新娘}
\verseline{}{已为你备好。别再为心爱的克瑞乌萨流泪。}
\verseline{785}{我不会看见密耳弥冬人或多洛佩斯人傲慢的居所}
\verseline{}{也不会前去侍奉希腊妇人——}
\verseline{}{我出自达耳达诺斯，又是女神维纳斯的儿媳。}
\verseline{}{是众神伟大的母亲把我留在这些海岸。}
\verseline{}{现在永别了，别忘了爱护我们共同的儿子。’}
\verseline{790}{她说完便舍下泪流不止、千言万语未出口的我，}
\verseline{}{退去，消散进轻薄的空气。}
\verseline{}{三次我试着用双臂环住她的颈项：\footnote{埃涅阿斯三次想抱住克瑞乌萨、三次扑空，沿用了《奥德赛》第十一卷奥德修斯三次拥抱亡母安提克勒娅的场面；第六卷他与安喀塞斯分别时又几乎逐字重复这三行。克瑞乌萨此前的“现在永别了”也让人想到维吉尔《农事诗》第四卷欧律狄刻与俄耳甫斯的诀别。}}
\verseline{}{三次，那徒然被抱住的幻影都逃出双手，}
\verseline{}{宛如轻风，又极似飞掠的梦境。}
\verseline{795}{就这样，黑夜耗尽后我终于回到同伴那里。}
\verseline{}{我惊讶地发现，大批新的同行者已经}
\verseline{}{涌来：有母亲、有男人，也有聚集起来}
\verseline{}{准备流亡的青年——一群可悲的人。}
\verseline{}{他们从四方汇聚，带着决心和资财，}
\verseline{800}{准备随我渡海，去往我愿意带他们抵达的任何土地。}
\verseline{}{晨星已从伊达山最高的山脊升起}
\verseline{}{带来白昼。达那俄斯人已控制各处城门，}
\verseline{}{再没有任何援救的希望。}
\verseline{}{我终于退让，背起父亲走向群山。”}
